%%This is a very basic article template.
%%There is just one section and two subsections.
%--------------------------------------------------------------
% thesis.tex 
%--------------------------------------------------------------
% Corso di Laurea in Informatica 
% http://if.dsi.unifi.it/
% @Facolt\`a di Scienze Matematiche, Fisiche e Naturali
% @Universit\`a degli Studi di Firenze
%--------------------------------------------------------------
% - template for the main file of Informatica@Unifi Thesis 
% - based on Classic Thesis Style Copyright (C) 2008 
%   Andr\'e Miede http://www.miede.de   
%--------------------------------------------------------------
\documentclass[oneside,titlepage,fleqn,
	headinclude,12pt,a4paper,BCOR5mm,footinclude]{scrbook}
%--------------------------------------------------------------
\newcommand{\myItalianTitle}{Principal component analysis: dalla teoria alla pratica\xspace} 
\newcommand{\myEnglishTitle}{\xspace}
%use the right myDegree option
\newcommand{\myDegree}{Corso di Laurea Magistrale in Informatica\\curriculum "Resilient and secure cyberphysical systems"\xspace}
\newcommand{\myName}{Marco Buracchi\xspace}
\newcommand{\myProf}{Prof.ssa Anna Gottard\xspace}
\newcommand{\myFaculty}{
	Scuola di Scienze Matematiche, Fisiche e Naturali\xspace}
\newcommand{\myUni}{\protect{
	Universit� degli Studi di Firenze}\xspace}
\newcommand{\myLocation}{Firenze\xspace}
\newcommand{\myTime}{Anno Accademico 2017-2018\xspace} 
\newcommand{\myVersion}{Version 0.1\xspace}

%--------------------------------------------------------------
\usepackage{dia-classicthesis-ldpkg}
%--------------------------------------------------------------
% Options for classicthesis.sty:
% tocaligned eulerchapternumbers drafting linedheaders 
% listsseparated subfig nochapters beramono eulermath parts 
% minionpro pdfspacing
\usepackage[eulerchapternumbers,linedheaders,subfig,beramono,eulermath,
parts]{classicthesis}
%--------------------------------------------------------------
\usepackage[italian]{babel}
\usepackage[latin1]{inputenc} 
\usepackage[T1]{fontenc} 
\usepackage[square,numbers]{natbib} 
\usepackage[fleqn]{amsmath}  
\usepackage{ellipsis}
\usepackage{subfig}
\usepackage{caption}
\usepackage{appendix}
\usepackage{siunitx}
\usepackage{stmaryrd}
\usepackage{tikz}
\usepackage[object=vectorian]{pgfornament}
\usepackage{listings}
%--------------------------------------------------------------
\newlength{\abcd} % for ab..z string length calculation
% how all the floats will be aligned
\newcommand{\myfloatalign}{\centering} 
\setlength{\extrarowheight}{3pt} % increase table row height
\captionsetup{format=hang,font=small}
%--------------------------------------------------------------
% Layout setting
%--------------------------------------------------------------
\usepackage{geometry}
\geometry{
	a4paper,
	ignoremp,
	bindingoffset = 1cm, 
	textwidth     = 13.5cm,
	textheight    = 21.5cm,
	lmargin       = 3.5cm, % left margin
	tmargin       = 4cm    % top margin 
}

\lstset{
  	frame=tb,
	language=Matlab,
  	aboveskip=3mm,
  	belowskip=3mm,
  	showstringspaces=false,
  	columns=flexible,
  	basicstyle={\small\ttfamily},
  	numbers=none,
  	breaklines=true,
  	breakatwhitespace=true,
  	tabsize=3
}

\newcommand{\sectionline}{
	\begin{center}
		\resizebox{0.5\linewidth}{1ex}{
			\begin{tikzpicture}
				\node  (C) at (0,0) {};
				\node (D) at (9,0) {};
				\path (C) to [ornament=83] (D);
			\end{tikzpicture}
		}
	\end{center}
}
%--------------------------------------------------------------
\begin{document}
\frenchspacing
\raggedbottom
\pagenumbering{roman}
\pagestyle{plain}
%--------------------------------------------------------------
% Frontmatter
%--------------------------------------------------------------
%--------------------------------------------------------------
% titlepage.tex (use thesis.tex as main file)
%--------------------------------------------------------------
\begin{titlepage}
	\begin{center}
   	\large
      \hfill
      \vfill
      \begingroup
         \includegraphics[scale=0.15]{logo/LOGO}\\
			\spacedallcaps{\myUni} \\ 
			\myFaculty \\
			\vspace{0.5cm}
			\myDegree \\ 
			\vspace{0.5cm}
         \vspace{0.5cm}    
         MULTIVARIATE ANALYSIS AND STATISTICAL LEARNING    
      \endgroup 
      \vfill 
      \begingroup
      	\color{Maroon}\spacedallcaps{\myItalianTitle} \\ $\ $\\
      	\spacedallcaps{\myEnglishTitle}	
	\bigskip
      \endgroup
      \spacedlowsmallcaps{\myName}
      \vfill 
      \vfill
      \emph{\myProf}
      %Relatore: \emph{Andrea Bondavalli}\\
      %Correlatore: \emph{Andrea Ceccarelli}\\
      \vfill
      \vfill
      \myTime
      \vfill                      
	\end{center}        
\end{titlepage}   
%--------------------------------------------------------------
% back titlepage
%--------------------------------------------------------------
%   \newpage
%	\thispagestyle{empty}
%	\hfill
%	\vfill
%	\noindent\myName: 
%	\textit{\myItalianTitle,} 
%	\myDegree, 
%	%\textcopyright\ 
%	\myTime
%--------------------------------------------------------------
% back titlepage end
%--------------------------------------------------------------
\tableofcontents
\pagenumbering{arabic}
\setcounter{secnumdepth}{0}

\chapter{Principal Component Analysis}
	
	\section{Storia}
		svolgimento
\begin{appendices}
	\chapter{Appendice}
	\lipsum
\end{appendices}
\pagestyle{scrheadings}
%--------------------------------------------------------------
% Mainmatter
%--------------------------------------------------------------

		
%--------------------------------------------------------------
\end{document}
%--------------------------------------------------------------
